\documentclass[11pt, letterpaper]{article}
\usepackage[margin=0.9in]{geometry}
\usepackage{hyperref}
\usepackage{enumitem}
\usepackage{xcolor}
\usepackage{titlesec}
\usepackage{fancyhdr}
\usepackage{tcolorbox}

\definecolor{accent}{HTML}{1A5490}
\definecolor{highlightbg}{HTML}{FFF8DC}
\hypersetup{colorlinks=true, linkcolor=accent, urlcolor=accent}

\pagestyle{fancy}
\fancyhf{}
\rhead{STA 561D | Duke}
\lhead{Infinite Connections --- Executive Summary}
\rfoot{\thepage}
\renewcommand{\headrulewidth}{0.4pt}

\titleformat{\section}{\large\bfseries\color{accent}}{}{0em}{}
\titlespacing\section{0pt}{6pt}{4pt}

\newtcolorbox{keybox}{
    colback=highlightbg, colframe=yellow!60!orange,
    arc=2mm, boxrule=0.5pt,
    left=6pt, right=6pt, top=4pt, bottom=4pt,
}

\title{\textbf{Infinite Connections}\\[0.2em]
\large Generating New NYT-Style Word Puzzles at Scale}
\author{Abuzar \quad Adreama \quad Burak \quad Hengkai \quad Kaiwen \\[0.3em]
STA 561D --- Duke University}
\date{April 2026}

\begin{document}
\maketitle
\thispagestyle{fancy}
\vspace{-2em}

\section*{The problem}

The New York Times Connections puzzle shows the player 16 words and asks them
to sort the words into 4 groups of 4 based on a hidden theme. The themes can
be simple (four synonyms for ``hello''), tricky (four things that fit before the
word ``fish''), or outright playful (four Taylor Swift album titles). For this
project we built a system that \emph{creates} these puzzles automatically, at
scale, with a web app that lets anyone play them in the browser.

Generating a good Connections puzzle is harder than it looks. The 16 words must
fit together in exactly one way --- no ambiguity. The themes must feel fresh
but recognizable. And we cannot accidentally reproduce a puzzle the NYT has
already published, which the grading rubric treats as an automatic failure.

\section*{What we built --- five complementary methods}

Rather than pick one approach and hope it works, each teammate built a different
method. Together they cover the full design space from ``ask the AI everything''
to ``use no AI at all.''

\begin{itemize}[leftmargin=*, itemsep=0.25em]
\item \textbf{Adreama --- Iterative AI generation.} Asks the AI to build one
group at a time, remembering what it has already used. Built up a library of
\textbf{400+} puzzles with memory that persists across sessions, so duplicates
never creep in even across many runs.

\item \textbf{Burak --- Pattern-driven dictionary builder.} Uses no AI for
generation. Instead, pulls words from language databases (WordNet, ConceptNet,
Datamuse) and fills puzzle slots using explicit patterns (synonyms,
category membership, shared actions). Specializes in the green and blue
difficulty levels with thousands of candidate groups.

\item \textbf{Hengkai --- Quality-first yellow specialist.} Focused on the
easiest (yellow) difficulty, where cohesion and obviousness matter most. Uses
cached word embeddings and strict filters that throw out any word too obscure
or technical.

\item \textbf{Kaiwen --- Two pipelines, compared head to head.}
\emph{Pipeline A} follows the published ``Making New Connections'' paper: five
AI calls per puzzle. \emph{Pipeline B} is our new method, \textbf{Category-First
Retrieval (CFR)}: the AI proposes only the category names (one call per puzzle,
or zero if we reuse NYT categories), then math finds the matching words from an
expanded 14{,}877-word dictionary.

\item \textbf{Abuzar --- Multi-agent critic system.} Builds puzzles with one AI
agent and has a second agent critique the result. Good ideas survive, bad ones
get rewritten. Pairs this with a playable web server and a library of ``concept
inspirations'' to keep themes fresh.
\end{itemize}

Individual generators use different duplicate-prevention strategies:
Kaiwen and Hengkai check each candidate against the 554 past NYT puzzles
directly; the other three guard against \emph{self}-repetition via persistent
memory files but don't check the NYT corpus. As a final gate, every puzzle in
our \textbf{merged submitted dataset} is passed through Kaiwen's shared
validator, which flags any 16-word set that shares more than six words with a
past NYT puzzle \emph{and} any 4-word group that exactly matches a past NYT
group. That merged gate is what guarantees our submitted dataset cannot
reproduce a past puzzle.

\section*{Results}

\begin{keybox}
Across all methods we generated roughly \textbf{2{,}000+ validated puzzles}.
The best pipeline (CFR) achieves a \textbf{99\% pass rate} on automated
validation, reproduces \textbf{zero} past NYT puzzles, and uses words from
a \textbf{14{,}877-word vocabulary --- three times larger than the NYT
original}. At full scale, generating 10{,}000 puzzles costs between
\textbf{\$0 and \$50} depending on the method, so we can produce as many as
the course requires.
\end{keybox}

\section*{What we learned}

The surprising finding was that the simplest, least-AI-heavy method produced
the highest pass rate. Asking the AI to do everything (pick words, pick
categories, pick difficulty) is both expensive and unreliable: the AI gives
inconsistent word choices and needs a separate editor pass to clean up. But
letting a language AI do only \emph{the creative naming step} --- and letting
math handle the mechanical word-lookup step --- produces better results at
one-twentieth the cost and four times the speed.

Across the team, the strongest puzzles tended to come from methods that
combined an AI's creativity with deterministic math for selection: everyone's
pipeline used embeddings (a way to compare how similar two words are
numerically) somewhere in the process.

\section*{Deliverables}

\begin{itemize}[leftmargin=*, itemsep=0.15em]
\item A \textbf{live web app} at
\url{https://kevin2330.github.io/nyt-connections-generator/}
loads 540+ validated puzzles in the NYT Connections interface. Anyone can play
them in a browser.
\item A \textbf{master Jupyter notebook} (\texttt{notebooks/master\_demo.ipynb})
walks through every method end-to-end in reproducible form. It runs without
any API keys --- all AI calls fall back to realistic mock responses when none
are configured.
\item A \textbf{technical report} (see Methods document and the per-teammate
folders) details the algorithms, math, and benchmarks.
\item A \textbf{large pre-generated dataset} on disk, from which the course can
sample puzzles for TA evaluation.
\end{itemize}

\section*{Future work}

The next step is \emph{human} evaluation, not just the automated solvers: put
the puzzles in front of real puzzle players and measure how many they would
believe came from the NYT. We added a mechanism for that to the web app
(thumbs-up / thumbs-down on each puzzle) but haven't yet collected a large
sample. A stretch goal is a fine-tuned embedding model trained specifically on
Connections puzzles, which would likely push the pass rate even higher for
puzzles built around wordplay (the one category where our current system is
weakest).

\end{document}
